\section{Planned Measures}
\label{sec:measures}

\subsection{Primary Outcome}

The proposed primary outcome is a multi-item perceived-fairness scale
administered after each map. Candidate items ask whether boundaries appear
reasonable, communities appear equally treated, the map appears manipulated,
the process inspires trust, and the districts reflect recognizable geography.

Item wording, response options, reverse coding, missing-item rules, scale
construction, and standardization will be finalized through cognitive testing
and a real pilot. No reliability, factor-analysis, or pilot estimate is
currently available.

\subsection{Secondary Outcomes}

Planned secondary outcomes include:

\begin{itemize}
\item perceived partisan advantage;
\item process trust;
\item willingness to accept an outcome that disadvantages the respondent's
party;
\item support for using a benchmark within a commission or court process; and
\item understanding of the benchmark's limits.
\end{itemize}

The last measure is necessary because a treatment that increases trust by
causing respondents to believe an algorithm guarantees fairness would be a
misleading communication effect rather than evidence of informed legitimacy.

\subsection{Moderators}

Preregistered moderators may include party identification, in-party status,
political knowledge, education, state familiarity, baseline institutional
trust, and algorithm attitudes. Subgroup estimates will be labeled exploratory
unless the study is powered and corrected for the full moderator family.

\subsection{Manipulation And Balance Checks}

Manipulation checks will test whether respondents understood the stated map
provenance and the benchmark's actual claim boundary. Balance tables and
exclusion summaries will be generated from observed data.

No manipulation success rate, randomization balance result, Cronbach's alpha,
or robustness result is reported in this protocol.
